\section*{Chapitre \ref{ch:b2:biogeochemical-cycles} --- Les cycles biogéochimiques}

\begin{solution}{exo:b2:biogeochemical-cycles:1}
Un \omterm{def:b2:biogeochemical-cycles:reservoir}{réservoir} est un stock de l'élément (une masse)~; un \omterm{def:b2:biogeochemical-cycles:reservoir}{flux} est une
vitesse de transfert entre \omterm{def:b2:biogeochemical-cycles:reservoir}{réservoirs} (une masse par an)~; le \omterm{def:b2:biogeochemical-cycles:reservoir}{temps de
résidence} est la masse du \omterm{def:b2:biogeochemical-cycles:reservoir}{réservoir} divisée par son \omterm{def:b2:biogeochemical-cycles:reservoir}{flux} sortant total,
la durée moyenne qu'y passe un atome. Plantes terrestres~: $450/120 \approx 4$ ans
(à peu près le carbone d'une feuille~; le \omterm{def:b2:plant-meristems:wood}{bois} le garde des décennies et
la moyenne masque la dispersion). Océan profond~: $37000/100 \approx 400$ ans
--- le temps qu'il faut à une molécule qui a sédimenté pour revoir la
surface.
\end{solution}

\begin{solution}{exo:b2:biogeochemical-cycles:2}
$2.5\times 2.13 = \qty{5.3}{GtC}$ par an. $10/2.13 = \qty{4.7}{ppm}$
--- si tout restait dans l'air~; avec une \omterm{prop:b2:biogeochemical-cycles:carbon}{fraction atmosphérique} de
\qty{45}{\%}, \qty{2.1}{ppm}.
\end{solution}

\begin{solution}{exo:b2:biogeochemical-cycles:3}
Ouverture~: la fixation biologique par la nitrogénase (cyanobactéries,
rhizobiums, \emph{Azotobacter}, \emph{Frankia}) et la foudre. Retour~:
la \omterm{def:b2:microbial-metabolism:anaerobic}{dénitrification} (anaérobies facultatifs comme \emph{Pseudomonas}
et \emph{Paracoccus} dans les sols et sédiments anoxiques) et l'anammox
(planctomycètes). Entre les deux~: ammonification de l'azote organique en
ammonium par les décomposeurs (champignons, bactéries)~; nitrification
de l'ammonium en nitrite (\emph{Nitrosomonas}, archées oxydant
l'ammoniac) et du nitrite en nitrate (\emph{Nitrobacter})~; et
assimilation de l'ammonium ou du nitrate par les plantes et les microbes.
\end{solution}

\begin{solution}{exo:b2:biogeochemical-cycles:4}
À l'échelle des temps géologiques, l'azote peut toujours être tiré de
l'inépuisable $\mathrm{N_2}$ de l'air par les fixateurs, qui sont favorisés
dès que l'azote manque~; le phosphore n'a pas de tel \omterm{def:b2:biogeochemical-cycles:reservoir}{réservoir} et son
apport est fixé par l'altération, si bien que le plafond à long terme de
la productivité est le phosphore. Au cours d'une saison, la fixation est
lente et coûteuse, et la plupart des plantes en sont incapables~: c'est
donc l'azote disponible qui limite la croissance ici et maintenant ---
c'est pourquoi les engrais sont surtout azotés.
\end{solution}

\begin{solution}{exo:b2:biogeochemical-cycles:5}
$\tau = 1/0.2 = 5$ ans~; $M^{*} = F\tau = \qty{10}{t}$, soit
$\qty{1}{g/m^3} = \qty{1000}{mg/m^3}$, fortement eutrophe. Temps pour
atteindre \qty{90}{\%}~: $\tau\ln 10 = 11.5$ ans. Diviser l'apport par
deux divise l'\omterm{def:b2:biogeochemical-cycles:reservoir}{état stationnaire} par deux, à \qty{500}{mg/m^3}, atteint
sur la même échelle de temps.
\end{solution}

\begin{solution}{exo:b2:biogeochemical-cycles:6}
Années 1960~: \qty{0.9}{ppm/yr}, \qty{1.9}{GtC/yr}. Années 2010~:
\qty{2.3}{ppm/yr}, \qty{5.0}{GtC/yr}. Le taux de croissance a été
multiplié par deux et demi, au rythme des émissions.
\end{solution}

\begin{solution}{exo:b2:biogeochemical-cycles:7}
\qty{6}{ppm} représentent \qty{13}{GtC}, contre une production primaire
nette de \qty{35}{GtC}~: les dents de scie en représentent un tiers. Elles
sont plus petites parce que la respiration et la décomposition se
poursuivent pendant l'été (les dents de scie enregistrent l'absorption
nette, non brute), parce que les saisons de l'hémisphère sud sont
inversées et compensent en partie, et parce que l'océan amortit
l'oscillation.
\end{solution}

\begin{solution}{exo:b2:biogeochemical-cycles:8}
$\qty{200}{kg}/(\qty{28}{g/mol}) = 7100$~mol de $\mathrm{N_2}$~; $16\times
7100 = \num{114000}$~mol d'ATP~; $/30 = 3800$~mol de glucose, \qty{690}{kg}.
Une culture qui produit \qty{10}{t} de matière sèche fixe à peu près
\qtyrange{15}{20}{t} de sucre une fois la respiration comptée, si bien
que la fixation lui coûte quelque \qty{4}{\%} de ses photosynthétats ---
le prix de son indépendance vis-à-vis de l'azote du sol.
\end{solution}

\begin{solution}{exo:b2:biogeochemical-cycles:9}
Rapport disponible $\mathrm{N:P} = 32/2.2 = 14.5$, inférieur aux 16 de
Redfield~: le nitrate s'épuise le premier, laissant
\qty{0.2}{\micro mol/kg} de phosphate. Carbone fixé~: $32\times 106/16 =
\qty{212}{\micro mol/kg}$, environ \qty{2.5}{mg} de carbone par
kilogramme d'eau.
\end{solution}

\begin{solution}{exo:b2:biogeochemical-cycles:10}
$E^{*} = Q\tau = \qty{500}{GtC}$, soit \qty{235}{ppm} au-dessus des 285
préindustriels~: environ \qty{520}{ppm}. La réalité est pire parce que les
puits ne forment pas un compartiment rapide unique~: l'océan superficiel
sature (sa capacité d'absorber le $\mathrm{CO_2}$ diminue à mesure que son
carbonate est consommé), le puits terrestre s'affaiblit avec le
réchauffement, et l'élimination réelle a une composante lente --- le
mélange avec l'océan profond et l'altération des roches --- avec des
constantes de temps de mille à cent mille ans, si bien qu'une fraction
de l'excédent ne se résorbe sur aucun $\tau$ de quelques décennies.
\end{solution}

\begin{solution}{exo:b2:biogeochemical-cycles:11}
Ce qui a été mesuré est l'échange du $\mathrm{CO_2}$ atmosphérique avec
l'océan et la biosphère, et non la désintégration~: le $\mathrm{{}^{14}C}$ des
bombes a été dilué dans les compartiments bien plus grands de l'océan
superficiel, des plantes et des sols, qui ont restitué du carbone
ordinaire. La demi-vie de 11 ans de l'excédent (un temps caractéristique
d'environ 16 ans) est l'échelle de temps de cet échange --- le
renouvellement du carbone de l'atmosphère vis-à-vis de l'océan
superficiel et des terres, plus long que le \omterm{def:b2:biogeochemical-cycles:reservoir}{temps de résidence} brut de
quatre ans parce qu'une grande partie de ce qui part revient aussitôt.
\end{solution}

\begin{solution}{exo:b2:biogeochemical-cycles:12}
Azote~: fixation naturelle d'environ 200~Tg par an~; fixation humaine
(Haber--Bosch 120, cultures 60, combustion 30) d'environ 210~: le \omterm{def:b2:biogeochemical-cycles:reservoir}{flux} est
doublé, et les \omterm{def:b2:biogeochemical-cycles:reservoir}{réservoirs} touchés --- sols, rivières, mers côtières, le
$\mathrm{N_2O}$ de l'air --- sont petits et rapides, si bien que le changement
apparaît en quelques décennies, tandis que le \omterm{def:b2:biogeochemical-cycles:reservoir}{réservoir} de $\mathrm{N_2}$ reste
intact. Carbone~: les émissions humaines de 11~GtC par an ne représentent
que \qty{5}{\%} de l'échange naturel brut d'environ 200, si bien que le
\omterm{def:b2:biogeochemical-cycles:reservoir}{flux} n'est qu'effleuré~; mais les \omterm{def:b2:biogeochemical-cycles:reservoir}{flux} naturels se compensent et le \omterm{def:b2:biogeochemical-cycles:reservoir}{flux}
humain non, et il a accru le \omterm{def:b2:biogeochemical-cycles:reservoir}{réservoir} atmosphérique de \qty{50}{\%} ---
la petite poussée a un effet cumulé important sur un \omterm{def:b2:biogeochemical-cycles:reservoir}{réservoir} dont le
puits ultime est lent. Savoir quel cycle est «~le plus perturbé~» dépend
de ce que l'on compte~: les \omterm{def:b2:biogeochemical-cycles:reservoir}{flux} ou l'accumulation.
\end{solution}

\begin{solution}{pb:b2:biogeochemical-cycles:1}
\textbf{1.} $285\times 2.13 = \qty{607}{GtC}$~; $412\times 2.13 =
\qty{878}{GtC}$~; augmentation \qty{271}{GtC}.
\textbf{2.} $271/700 = 0.39$ (la fraction au cours des dernières
décennies, émissions liées à l'usage des terres comprises, est plus
proche de $0.45$).
\textbf{3.} L'océan, environ un quart des émissions, et les terres
(repousse des forêts et fertilisation par le $\mathrm{CO_2}$), environ un
tiers~; ensemble, les \qty{430}{GtC} qui ne sont pas dans l'air.
\textbf{4.} $175/38000 = \qty{0.5}{\%}$ --- négligeable pour l'océan
entier~; mais l'absorption s'est jusqu'ici faite surtout dans la couche
de surface (\qty{900}{GtC}), où elle représente une hausse d'un dixième
ou plus, et la chimie des carbonates transforme une petite hausse du
$\mathrm{CO_2}$ dissous en une hausse plus grande de la concentration en ions
hydrogène~: le pH de surface a baissé de $0.1$, soit une hausse de
\qty{30}{\%} de l'acidité.
\textbf{5.} Avec $Q = Q_0 e^{t/T}$, essayer $E = Ae^{t/T}$~: $A/T = Q_0 -
A/\tau$, donc $A = Q_0\tau T/(\tau + T)$ et la \omterm{prop:b2:biogeochemical-cycles:carbon}{fraction atmosphérique}
$\dot E/Q = \tau/(\tau + T)$ est constante. Des émissions croissant de
\qty{2}{\%} par an ($T = 50$) avec $\tau = 40$ donnent $0.44$~: la
fraction est constante parce que le puits et la source croissent ensemble.
\textbf{6.} Avec $Q$ constant, $E \to Q\tau$ et $\dot E \to 0$~: la
\omterm{prop:b2:biogeochemical-cycles:carbon}{fraction atmosphérique} tend vers zéro à mesure que les puits rattrapent
leur retard --- dans le \omterm{thm:b2:biogeochemical-cycles:box}{modèle à un compartiment}~; en réalité, les puits
saturent et la fraction diminue plus lentement.
\textbf{7.} \qty{271}{GtC}.
\textbf{8.} $E^{*} = 10\times 100 = \qty{1000}{GtC}$~: $285 +
1000/2.13 = \qty{755}{ppm}$.
\textbf{9.} $E(t) = 1000 + (271 - 1000)e^{-t/100}$~; pour $t = 80$,
$E = 1000 - 729\times 0.449 = \qty{672}{GtC}$~: \qty{600}{ppm}.
\textbf{10.} $Q = 10(1 - t/50)$~: solution particulière $E_p = at + b$
avec $a = -20$, $b = 3000$ (d'après $a = 10 - b/\tau$ et $0 = -0.2 -
a/\tau$)~; $E = 3000 - 20t + (271 - 3000)e^{-t/100}$. Pour $t = 50$~:
$2000 - 2729\times 0.607 = \qty{344}{GtC}$, \qty{447}{ppm}. Ensuite, $E$
décroît librement~: pour $t = 80$, $344\,e^{-0.3} = \qty{255}{GtC}$,
\qty{405}{ppm}.
\textbf{11.} \qty{600}{ppm} contre \qty{405}{ppm}~: le second scénario
ramène l'air, en 2100, presque à son niveau actuel.
\textbf{12.} Le $\tau$ unique fait travailler les puits au même rythme
sur un excédent qui diminue~; en réalité, les puits rapides (la couche
de mélange, les forêts en croissance) ont déjà pris ce qu'ils pouvaient,
et l'élimination restante se fait par le lent mélange avec l'océan
profond et par l'altération, avec des $\tau$ de plusieurs siècles à
plusieurs millénaires, si bien que la décroissance après la neutralité
carbone est bien plus lente --- la concentration reste à peu près
constante pendant des siècles au lieu de baisser.
\textbf{13.} Culture \qty{90}{kg}~; dénitrifié \qty{37.5}{kg}~; lessivé
\qty{22.5}{kg} d'azote par hectare et par an.
\textbf{14.} $\qty{22.5}{kg}/\qty{14}{g/mol} = 1600$~mol de N~; carbone
$1600\times 106/16 = \num{10600}$~mol, \qty{128}{kg}.
\textbf{15.} \num{10600}~mol de $\mathrm{O_2}$, \qty{340}{kg}.
\textbf{16.} $\num{10600}/0.25 = \num{42000}$~\unit{m^3}~: le lessivage
d'un hectare peut priver d'oxygène quarante mille mètres cubes d'eau de
fond chaque année --- et le Mississippi draine cent millions d'hectares.
\textbf{17.} $37.5\times 0.01 = \qty{0.375}{kg}$ d'azote sous forme de
$\mathrm{N_2O}$, soit $0.375\times 44/28 = \qty{0.59}{kg}$ de $\mathrm{N_2O}$~;
$\times 270 = \qty{160}{kg}$ d'équivalent $\mathrm{CO_2}$.
\textbf{18.} $150\times 4 = \qty{600}{kg}$ de $\mathrm{CO_2}$~: la
fabrication pèse quatre fois plus que le $\mathrm{N_2O}$ du champ, et les
deux ensemble représentent trois quarts de tonne par hectare et par an.
\textbf{19.} $\tau = 10$ ans~; $M^{*} = 5\times 10 = \qty{50}{t}$~;
concentration $50/(5\times 10^{7}) = \qty{1}{g/m^3} =
\qty{1000}{mg/m^3}$.
\textbf{20.} Trente fois le seuil~: fortement eutrophe.
\textbf{21.} $1/\tau = 0.1 + 0.2 = 0.3$~: $\tau = 3.3$ ans~; $M^{*} =
5/0.3 = \qty{16.7}{t}$, \qty{330}{mg/m^3}.
\textbf{22.} $M^{*} = 1/0.3 = \qty{3.3}{t}$, \qty{67}{mg/m^3} --- encore
eutrophe~; à moins de \qty{10}{\%} au bout de $\tau\ln 10 = 7.7$ ans.
\textbf{23.} Le sédiment n'est pas un puits à sens unique~: le phosphore
qui s'y est accumulé pendant les années d'eutrophisation est relargué
dans les conditions anoxiques que crée l'eutrophisation elle-même (le
phosphate lié au fer se dissout lorsque le fer est réduit), si bien que
le lac s'alimente lui-même pendant des années après la suppression de
l'apport extérieur --- la charge interne, un second \omterm{def:b2:biogeochemical-cycles:reservoir}{réservoir} au long
\omterm{def:b2:biogeochemical-cycles:reservoir}{temps de résidence}.
\textbf{24.} $\qty{5000}{kg}/\qty{31}{g/mol} = 1.6\times 10^{5}$~mol de P~;
carbone $\times 106 = 1.7\times 10^{7}$~mol, \qty{205}{t} de carbone par
an --- \qty{20}{g/m^2} sur la surface du lac s'il est profond de
\qty{5}{m}, une forte prolifération.
\textbf{25.} \omterm{prop:b2:biogeochemical-cycles:carbon}{Fraction atmosphérique} $0.39$ sur 1850--2020~; en 2100,
\qty{600}{ppm} avec des émissions maintenues à \qty{10}{GtC/yr} et
\qty{405}{ppm} avec des émissions ramenées à zéro en 2070 (un modèle
optimiste à $\tau$ unique)~; azote lessivé \qty{22.5}{kg} par hectare
et par an.
\end{solution}
